% from ega4-8-fr.tex
\oldpage[IV-3]{40}
un $A$-morphisme séparé. Alors il existe deux $A$-préschémas $X'$, $P$ de
présentation finie, et des $A$-morphismes $p:P\to Y$, $j:X'\to P$,
$g:X'\to X$, tels que le diagramme
\[
\begin{tikzcd}[column sep=large,row sep=large]
X' \arrow[r,"j"] \arrow[d,"g"'] & P \arrow[d,"p"] \\
X \arrow[r,"f"'] & Y
\end{tikzcd}
\]
soit commutatif, et que~: 1\textsuperscript{o} $p$ soit projectif~;
2\textsuperscript{o} $g$ soit projectif et surjectif~;
3\textsuperscript{o} $j$ soit une immersion ouverte.
\end{proposition}

\begin{proof}
En effet, soient $A_0\subset A$, $X_0$, $Y_0$ et $f_0$ déterminés comme dans
(8.9.1) de sorte que $Y_0$ est noethérien et $f_0$ de type fini~; on peut en
outre supposer $f_0$ séparé d'après (8.10.4). Le lemme de Chow (II, 5.6.1)
montre alors l'existence de trois morphismes $p_0:P_0\to Y_0$,
$g_0:X'_0\to X_0$ et $j_0:X'_0\to P_0$, de type fini, tels que le diagramme
\[
\begin{tikzcd}[column sep=large,row sep=large]
X'_0 \arrow[r,"j_0"] \arrow[d,"g_0"'] & P_0 \arrow[d,"p_0"] \\
X_0 \arrow[r,"f_0"'] & Y_0
\end{tikzcd}
\]
soit commutatif, et que $p_0$ soit projectif, $g_0$ projectif et surjectif et
$j_0$ une immersion ouverte. Les propriétés de l'énoncé résultent alors de
l'invariance des propriétés précédentes par changement de base (II, 5.5.5,
(iii) et I, 3.5.2 et 4.3.2).
\end{proof}

\begin{proof}[Suite de la démonstration de (8.10.5)]
\emph{Cas (xii)}~: Appliquons au morphisme $f_0:X_0\to S_0$ la prop.
(8.10.5.1)~: on a donc un diagramme commutatif
\[
\begin{tikzcd}[column sep=large,row sep=large]
X'_0 \arrow[r,"j_0"] \arrow[d,"g_0"'] & P_0 \arrow[d,"p_0"] \\
X_0 \arrow[r,"f_0"'] & S_0
\end{tikzcd}
\]
où $p_0$ est projectif, $g_0$ projectif et surjectif, et $j_0$ une immersion
ouverte~; on en déduit pour chaque $\lambda$ un diagramme analogue où les
morphismes $p_\lambda=p_0\times1_{S_\lambda}$,
$g_\lambda=g_0\times1_{S_\lambda}$ et $j_\lambda=j_0\times1_{S_\lambda}$ ont
respectivement les mêmes propriétés, et de même pour les morphismes limites
projectives $p=p_0\times1_S$, $g=g_0\times1_S$, $j=j_0\times1_S$. Comme $g$
est propre (II, 5.5.3), il en est de même de $p\circ j=f\circ g$ (II, 5.4.2)
et comme $p$ est séparé, $j$ est propre, donc une immersion fermée~;
appliquant le cas (iv) au morphisme $j$ (en notant que $X'_0$ et $P_0$ sont
des $S_0$-préschémas de présentation finie (8.10.5.1 et 1.6.2)), on voit
qu'il existe $\lambda$ tel que $j_\lambda$ soit une immersion fermée, donc
soit propre (II, 5.4.2). Mais alors
$f_\lambda\circ g_\lambda=j_\lambda\circ p_\lambda$ est propre (II, 5.5.3 et
5.4.2) et comme $g_\lambda$ est surjectif, et qu'on peut supposer $f_\lambda$
séparé en vertu de l'hypothèse sur $f$ et de (v), il résulte de (II, 5.4.3)
que $f_\lambda$ est propre.

\emph{Cas (xiii) et (xiv)}~: En vertu de (xii) et de (II, 5.5.3) (qui est
applicable puisque les $S_\lambda$ sont quasi-compacts et quasi-séparés,
compte tenu de (1.7.19)), il suffit de
